% ================================================================
% SECTION 3: QUICK START & WORKFLOW
% ================================================================
\section{Quick Start \& Workflow}
\label{sec:quickstart}

\subsection{Installation}

\bionetflux{} uses a standard \code{pyproject.toml} build.  Install in development (editable) mode:

\begin{lstlisting}[language=bash, caption={Installation}]
cd BioNetFlux
pip install -e .          # installs bionetflux + dependencies
pip install -e ".[dev]"   # also installs pytest, black, mypy, etc.
\end{lstlisting}

After installation the package \code{bionetflux} and the top-level shim \code{setup\_solver} are both importable from anywhere.

\subsection{Minimal Working Example}

A complete simulation requires four steps: \textbf{setup}, \textbf{time-stepper creation}, \textbf{time evolution loop}, and \textbf{visualisation}.

\begin{lstlisting}[language=Python, caption={Minimal working example}]
from setup_solver import quick_setup
from bionetflux.time_integration.time_stepper import TimeStepper
from bionetflux.visualization.lean_matplotlib_plotter import LeanMatplotlibPlotter

# 1. Setup (loads problem module, geometry, config)
setup = quick_setup(
    problem_module="bionetflux.problems.ks_problem",
    validate=True,
    config_file="config/ks_parameters.toml"
)

# 2. Create time stepper and initial conditions
stepper = TimeStepper(setup, verbose=True)
solution, bulk_data = stepper.initialize_solution()

# 3. Time loop
dt = setup.global_discretization.dt
T  = setup.global_discretization.T
t  = 0.0
while t < T:
    result = stepper.advance_time_step(solution, bulk_data, t, dt)
    if result.converged:
        solution  = result.updated_solution
        bulk_data = result.updated_bulk_data
        t += dt
    else:
        print(f"Newton failed at t={t:.4f}")
        break

# 4. Visualise
traces, multipliers = setup.extract_domain_solutions(solution)
plotter = LeanMatplotlibPlotter(
    setup.problems,
    setup.global_discretization.spatial_discretizations
)
plotter.plot_2d_curves(traces, title="Final solution")
plotter.show_all()
\end{lstlisting}

\subsection{TOML Configuration Overview}

Physical parameters, discretisation settings, initial/boundary conditions, and source terms are specified in TOML files.  A Keller-Segel configuration (\code{config/ks\_parameters.toml}) has the following top-level sections:

\begin{lstlisting}[caption={TOML configuration structure}]
[problem]                   # name, neq, problem_type
[time_parameters]           # T, dt
[discretization]            # n_elements or h, tau
[physical_parameters.diffusion]     # nu, mu
[physical_parameters.reaction]      # a, b
[physical_parameters.chemotaxis]    # chi, dchi (function names)
[initial_conditions]        # per-equation function names
[exact_solutions]           # per-equation function names (optional)
[exact_solution_derivatives]# per-equation function names (optional)
[force_functions]           # per-equation function names
[domain_initial_conditions] # domain-specific overrides
[domain_force_functions]    # domain-specific overrides
[boundary_conditions]       # per-boundary-point overrides
\end{lstlisting}

Function values can be:
\begin{itemize}
    \item A \textbf{named function} from the built-in library (e.g., \code{"zeros"}, \code{"ones"}, \code{"traveling\_wave\_u"}, \code{"gaussian"}).
    \item A \textbf{SymPy expression string} using variables \code{s} (space) and \code{t} (time), e.g., \code{"2.5 + 0.0*s"}.
    \item A \textbf{constant string} like \code{"constant"} (resolves to a constant-one function).
\end{itemize}

See Section~\ref{sec:config} for the full configuration reference.

\subsection{Adding a Custom Problem}

To create a new problem type:

\begin{enumerate}
    \item Create a Python module in \code{bionetflux/problems/} with a function:
    \begin{lstlisting}[language=Python]
def create_global_framework(geometry=None, config_file=None):
    # ... build problems, discretizations, constraints ...
    return problems, global_discretization, constraint_manager, "My Problem"
    \end{lstlisting}
    \item If the PDE system requires a new static condensation procedure, sub-class and register \code{StaticCondensationBase}:
    \begin{lstlisting}[language=Python]
from bionetflux.core.static_condensation_factory import StaticCondensationFactory
StaticCondensationFactory.register_implementation("my_type", MyStaticCondensation)
    \end{lstlisting}
    \item Call \code{quick\_setup("bionetflux.problems.my\_module")} from a driver script.
\end{enumerate}
